% Autonomous agents and systems cover a number of application areas, from robotics and digital assistants to combinatorial optimization, all sharing common, unresolved research challenges.
% It is not sufficient for agents to merely solve a given task; they must generalize to out-of-distribution tasks, perform reliably, and use hardware resources efficiently during training and on-device deployment, among other requirements.
% Several major classes of methods, such as reinforcement learning and imitation learning, are commonly used to tackle these problems, each with different trade-offs.
% However, there is currently no benchmarking suite that defines the environments, datasets, and metrics which can be used to develop reference implementations and seed leaderboards with baselines, providing a meaningful way for the community to compare progress.
% We introduce \method---a benchmarking suite including three environments that closely resemble real-world domains: computer chip floorplanning, web navigation, and quadruped locomotion.
% \method~provides metrics that track task performance, generalization, system resource efficiency, and reliability, which are all critical to real-world applications.
% In addition, we propose a data cost metric to account for the cost incurred acquiring offline data for imitation learning, reinforcement learning, and hybrid algorithms, which allows us to better compare these approaches.
% \method~also contains baseline implementations of standard algorithms, enabling apples-to-apples comparisons across methods and facilitating progress in real-world autonomy.
% As an open-source and extendable benchmark, \method~is designed to remain accessible, documented, up-to-date, and useful to the research community over the long term.

% Autonomous agents and systems cover a number of application areas, from robotics and digital assistants to combinatorial optimization, all sharing common, unresolved research challenges. It is not sufficient for agents to merely solve a given task; they must generalize to out-of-distribution tasks, perform reliably, and use hardware resources efficiently during training and on-device deployment, among other requirements.
% Several major classes of methods, such as reinforcement learning and imitation learning, are commonly used to tackle these problems, each with different trade-offs. However, there is currently no benchmarking suite that defines the environments, datasets, and metrics which can be used to develop reference implementations and seed leaderboards with baselines, providing a meaningful way for the community to compare progress.
% We introduce \method—a benchmarking suite including three environments that closely resemble real-world domains: computer chip floorplanning, web navigation, and quadruped locomotion.
% \method provides metrics that track task performance, generalization, system resource efficiency, and reliability, which are all critical to real-world applications.
% Using \method, we demonstrate that web navigation agents can achieve latencies comparable to human reaction times on consumer hardware,
% reveal important reliability trade-offs between algorithms for quadruped locomotion, and quantify the total energy costs of different learning approaches for computer chip-design.
% In addition, we propose a data cost metric to account for the cost incurred acquiring offline data for imitation learning, reinforcement learning, and hybrid algorithms, which allows us to better compare these approaches.
% \method also contains baseline implementations of standard algorithms, enabling apples-to-apples comparisons across methods and facilitating progress in real-world autonomy.
% As an open-source and extendable benchmark, \method is designed to remain accessible, documented, up-to-date, and useful to the research community over the long term.


Autonomous agents and systems cover a number of application areas, from robotics and digital assistants to combinatorial optimization, all sharing common, unresolved research challenges. It is not sufficient for agents to merely solve a given task; they must generalize to out-of-distribution tasks, perform reliably, and use hardware resources efficiently during training and on-device deployment, among other requirements.
Several classes of methods, such as reinforcement learning and imitation learning, are commonly used to tackle these problems, each with different trade-offs.
However, there is a lack of benchmarking suites that define the environments, datasets, and metrics which can be used to provide a meaningful way for the community to compare progress on applying these methods to real-world problems.
We introduce \method—a benchmarking suite including three environments that closely resemble real-world domains: computer chip floorplanning, web navigation, and quadruped locomotion.
\method provides metrics that track task performance, generalization, system resource efficiency, and reliability, which are all critical to real-world applications.
Using \method, we demonstrate that web navigation agents can achieve latencies comparable to human reaction times on consumer hardware,
reveal important reliability trade-offs between algorithms for quadruped locomotion, and quantify the total energy costs of different learning approaches for computer chip-design.
In addition, we propose a data cost metric to account for the cost incurred acquiring offline data for imitation learning, reinforcement learning, and hybrid algorithms, which allows us to better compare these approaches.
\method also contains baseline implementations of standard algorithms, enabling apples-to-apples comparisons across methods and facilitating progress in real-world autonomy.
As an open-source and extendable benchmark, \method is designed to remain accessible, documented, up-to-date, and useful to the research community over the long term.